\documentclass[11pt]{article}
\usepackage{fontspec}
\setmainfont{Times New Roman}
\setsansfont{Arial}
\setmonofont{Courier New}
\usepackage{amsmath,amssymb,amsthm,booktabs,longtable,array}
\usepackage{unicode-math}
\setmathfont{latinmodern-math.otf}
\usepackage[margin=1in]{geometry}
\usepackage{xcolor,hyperref,microtype}
\hypersetup{colorlinks=true,linkcolor=blue!45!black,urlcolor=blue!45!black}
\setlength{\emergencystretch}{3em}
\newtheorem{theorem}{Theorem}
\newtheorem{lemma}[theorem]{Lemma}
\newtheorem{corollary}[theorem]{Corollary}
\newcommand{\F}{\mathbb F}
\newcommand{\PG}{\operatorname{PG}}
\newcommand{\spanv}{\operatorname{span}}
\newcommand{\bad}{\mathbf 1_{\rm broken}}
\title{Rank-controlled projective-incidence repair\\by weighted chart transport}
\author{Robert Huynh}
\date{Working draft, 27 October 2026}
\begin{document}
\maketitle
\begin{abstract}
Let two arbitrary functions map the points and hyperplanes of canonical
$\PG(n,q)$ to those of $\PG(n,r)$, where $n\geq2$ and $q,r$ are prime
powers. Let $T$ count broken source flags, let $d$ be the source incidence
degree, and let $\rho$ be the fraction of ordered projective frames that
survive under the point map. If $\rho>0$ and the flag error fraction
$\delta$ is at most $\rho/(10^9n^2)$, both functions agree, apart from at
most $5T/q^{n-1}\leq10T/d$ entries in total, with one coupled field-semilinear
embedding. Neither input map is assumed injective. We give the weighted
plane correction, its transport to higher rank, and the paired incidence
gap used to obtain this estimate. The repair theorem and its two stated
formal corollaries have existing Lean verification. This is a separate
incidence-rigidity direction associated with the research on Erd\H{o}s
Problem 85; it does not settle that problem's eventual monotonicity question.
Public priority remains unresolved.
\end{abstract}
\tableofcontents

\section{Statement and scope}
Let $P$ and $\mathcal H$ be the point and hyperplane sets of canonical
$\PG(n,q)$. Throughout, source and target have the same projective dimension
$n\geq2$, while $q$ and $r$ are arbitrary prime powers. Put
\begin{equation}\label{counts}
 N=\frac{q^{n+1}-1}{q-1},\quad d=\frac{q^n-1}{q-1},\quad
 \lambda=\frac{q^{n-1}-1}{q-1},\quad a=d-\lambda=q^{n-1},\quad m=Nd.
\end{equation}
Consider arbitrary fixed functions
\[
 \phi:P\longrightarrow P(\PG(n,r)),\qquad
 \psi:\mathcal H\longrightarrow\mathcal H(\PG(n,r)).
\]
A source flag is a pair $(p,H)$ with $p\in H$. It is broken when
$\phi(p)\notin\psi(H)$. Let $T$ count these original flags once each and
put $\delta=T/m$. An ordered projective frame is an ordered list of $n+2$
points, every $n+1$ of which are independent. Let $\rho$ be the fraction
of uniform ordered source frames whose point images form target frames.
Repeated image points fail this condition.

For a field embedding $\sigma:\F_q\hookrightarrow\F_r$ and
$A\in\operatorname{GL}_{n+1}(\F_r)$, write $G$ for the point map
$[v]\mapsto[A\sigma(v)]$ together with its coupled hyperplane map. The
image of a hyperplane is the $\F_r$-linear span of the images of its
source points. Thus it need not equal their set. Define
\[
 D(G)=|\{p:\phi(p)\ne G(p)\}|+
       |\{H:\psi(H)\ne G(H)\}|.
\]
\begin{theorem}[Coupled repair]\label{main}
If $\rho>0$ and $\delta\leq\rho/(10^9n^2)$, a field embedding
$\F_q\hookrightarrow\F_r$ exists, and some coupled $G$ satisfies
\begin{equation}\label{bound}
 D(G)\leq\frac{5T}{q^{n-1}}\leq\frac{10T}{d}.
\end{equation}
\end{theorem}
\begin{corollary}[Collapse or repair]\label{dichotomy}
Every pair $(\phi,\psi)$ satisfies at least one alternative:
\[
 \text{some coupled }G\text{ has }D(G)\leq10T/d,
 \qquad\text{or}\qquad \rho\leq10^9n^2\delta.
\]
\end{corollary}
\begin{corollary}[Incompatible fields]\label{fields}
If there is no embedding $\F_q\hookrightarrow\F_r$, then
$\rho\leq10^9n^2\delta$ for every pair of maps.
\end{corollary}
The classical finite-field criterion makes the last hypothesis explicit:
the characteristics differ, or $q=p^u,r=p^v$ with $u\nmid v$.
The first corollary follows because failure of its second alternative
implies the hypotheses of Theorem~\ref{main}. The second then follows
because its repair alternative would supply a field embedding.

\section{A weighted protected-plane lemma}\label{plane}
First work in $\PG(2,q)$ and $\PG(2,r)$. Suppose the point and line maps
fix, in the corresponding source and target coordinates, the seven points
\[
\begin{gathered}
 U=[1:0:0],\quad V=[0:1:0],\quad O=[0:0:1],\quad I=[1:1:1],\\
 W=[1:1:0],\quad E=[1:0:1],\quad J=[0:1:1]
\end{gathered}
\]
and the six lines $\ell_\infty$, $d_0:y=x$, $v_0:x=0$, $v_1:x=1$,
$h_0:y=0$, $h_1:y=1$. Write $(x,y)$ for $[x:y:1]$, and set
\[
 v_x:\ X=x,\quad h_y:\ Y=y,\quad d_b:\ Y=X+b,\quad
 \ell_a:\ Y=aX,\quad\ell_{a,b}:\ Y=aX+b.
\]
Let $b(p,\ell)$ be the broken-flag indicator. The cost $\eta$ is the sum
of the 24 normalized sums listed in Table~\ref{families}. A family with
$k$ varying field parameters has denominator $q^k$. Restricted nonzero
parameters retain this denominator, rather than replacing $q$ by $q-1$.
Each displayed flag in a row is a separate sum. Repetitions are intentional.
\begin{table}[htbp]
\centering\small
\begin{tabular}{@{}clp{9.2cm}@{}}\toprule
Sums&Parameters&Flags\\\midrule
2&$x$; $y$&$(V,v_x)$; $(U,h_y)$\\
2&$x,y$&$((x,y),v_x)$; $((x,y),h_y)$\\
2&$t$&$((t,t),v_t)$; $((t,t),h_t)$\\
1&$t$&$((t,t),d_0)$\\
1&$a\ne0$&$(O,\ell_a)$\\
1&$a\ne0,t$&$((t,at),\ell_a)$\\
2&$a\ne0$&$((1,a),v_1)$; $((1,a),h_a)$\\
2&$a\ne0$&$((a^{-1},1),v_{a^{-1}})$; $((a^{-1},1),h_1)$\\
2&$a\ne0$&$((1,a),\ell_a)$; $((a^{-1},1),\ell_a)$\\
1&$b$&$(W,d_b)$\\
1&$x,b$&$((x,x+b),d_b)$\\
4&$b$; $i=0,1$&$((i,i+b),v_i)$; $((i,i+b),h_{i+b})$, for each fixed $i$\\
2&$b$; $i=0,1$&$((i,i+b),d_b)$, for each fixed $i$\\
1&$a,b,x$&$((x,ax+b),\ell_{a,b})$\\\bottomrule
\end{tabular}
\caption{The 24 weighted flag families.}\label{families}
\end{table}

Define $B_X=\{x:(V,v_x)\text{ is broken}\}$,
$X=\{x\notin B_X:\psi(v_x)\ne\ell_\infty\}$, and
$X_0=\F_q\smallsetminus(B_X\cup X)$. Define $B_Y,Y,Y_0$ analogously with
$(U,h_y)$. For $x\in X$, the image vertical has equation $X=f(x)$;
for $y\in Y$, the image horizontal has equation $Y=g(y)$. Extend $f,g$
by zero on all other inputs. The protected values imply
\begin{equation}\label{protectedvalues}
 0,1\in X\cap Y,\qquad f(0)=g(0)=0,\qquad f(1)=g(1)=1.
\end{equation}
\begin{lemma}[Weighted plane correction]\label{planelemma}
If $\eta\leq10^{-7}$, there is a field embedding
$\sigma:\F_q\hookrightarrow\F_r$ with
\[
\begin{gathered}
 q-|X|\leq17\eta q,\qquad q-|Y|\leq17\eta q,\\
 |\{x:f(x)\ne\sigma(x)\}|\leq288\eta q,\qquad
 |\{y:g(y)\ne\sigma(y)\}|\leq648\eta q.
\end{gathered}
\]
The affine point map equals $(x,y)\mapsto(\sigma(x),\sigma(y))$ outside
at most $971\eta q^2$ points. The conclusion also holds when the cost
uses any nonnegative pointwise majorant of $b$.
\end{lemma}

\subsection{Support estimates}
Let $Q$ be the affine points with a broken vertical or horizontal
coordinate flag. In this subsection only, write $x=|X|$, $y=|Y|$,
$b_X=|B_X|$, and $b_Y=|B_Y|$. Outside $Q$, intersection of the two image
coordinate lines gives image $(f(u),g(v))$ for $(u,v)\in X\times Y$,
image $V$ on $X\times Y_0$, and image $U$ on $X_0\times Y$.
Let $Q_{\rm diag}$ count diagonal points in $Q$, and let $e_{\rm diag}$
count broken affine point--$d_0$ flags. Since the target diagonal contains
neither $U$ nor $V$,
\begin{equation}\label{symdiff}
 |X\mathbin{\triangle} Y|\leq b_X+b_Y+Q_{\rm diag}+e_{\rm diag}\leq\eta q.
\end{equation}
The two diagonal coordinate families dominate the union count in this inequality.
For $a\ne0$ put
$r_a=|\{t\in X:at\in Y_0\}|$ and
$s_a=|\{t\in X_0:at\in Y\}|$.
Let $e_O$ count broken $(O,\ell_a)$ flags and $E_{\rm pen}$ broken
$((t,at),\ell_a)$ flags. If the $O$ flag is good, the image line contains
$O$ and cannot contain both $U,V$. Thus at least $\min(r_a,s_a)$ mixed
points lie in $Q$ or have a broken pencil-line flag. Each such point has
two nonzero coordinates and a unique nonzero slope. For a broken $O$
flag use the upper bound $q$. Hence
\begin{equation}\label{mixed}
 \sum_{a\ne0}\min(r_a,s_a)\leq qe_O+|Q|+E_{\rm pen}\leq\eta q^2.
\end{equation}
Moreover $|r_a-s_a|\leq|x-y|+b_X+b_Y\leq2\eta q$. It follows that
$\sum_{a\ne0}r_a\leq3\eta q^2$. Counting the unique slope for each pair
in $(X\smallsetminus\{0\})\times Y_0$ gives
\begin{equation}\label{productsupport}
 (x-1)(q-y-b_Y)\leq3\eta q^2,\qquad (x-1)(q-y)\leq4\eta q^2.
\end{equation}
For $a\in Y_0\smallsetminus\{0\}$ with $a^{-1}\in X_0$, use the two points
$(1,a)$ and $(a^{-1},1)$ on $\ell_a$. Their images are $V,U$ unless a
coordinate flag fails. If those and both point-line flags and the $O$
flag are good, one target line contains $O,U,V$, a contradiction.
Let $R_Q$ count appearances in $Q$ in these two families and $R_F$ their
broken $\ell_a$ flags. The table pays for each of these counts. All other
nonzero $a$ lie in $Y\cup B_Y$ or have inverse in $X\cup B_X$. Therefore
\begin{equation}\label{largesupport}
 x+y\geq q+1-(b_X+b_Y+e_O+R_Q+R_F)\geq q+1-\eta q.
\end{equation}
There is no reciprocal exception at $a=1$, by \eqref{protectedvalues}.
For $q=2$ that equation already gives $X=Y=\F_q$. For $q\geq3$,
\eqref{symdiff} and \eqref{largesupport}, with $\eta\leq1/12$, give
$x-1\geq(q-1-2\eta q)/2\geq q/4$. Equation~\eqref{productsupport}
then proves $q-y\leq16\eta q$ and $q-x\leq17\eta q$.
These are inequalities on integer counts with no rounding remainder.

\subsection{Additive correction and protected values}
The product identity $\phi(u,v)=(f(u),g(v))$ fails at at most
\[
 K\leq q(q-|X|)+q(q-|Y|)+|Q|\leq35\eta q^2
\]
affine points. Define $c(b)$ as the target intercept of $\psi(d_b)$ if
its $W$ flag is good and its image is a proper affine line; otherwise set
$c(b)=0$. In the proper case its slope is 1. In particular $c(0)=0$.
Let $R$ count failures of $g(u+b)=f(u)+c(b)$. Each product exception is
used once, a broken $W$ flag at most $q$ times, and a broken
point--$d_b$ flag once. Thus, writing $e_W,E_{\rm diag}$ for the latter
counts,
\begin{equation}\label{pexidererror}
 R\leq K+qe_W+E_{\rm diag}\leq36\eta q^2,\qquad r_0=R/q^2\leq36\eta.
\end{equation}
For each fixed $i=0,1$, the protected vertical line gives the sharper
row bound
\begin{equation}\label{fixedrow}
 |\{b:g(i+b)\ne f(i)+c(b)\}|\leq20\eta q.
\end{equation}
Indeed only missing horizontal coordinates, coordinate-flag failures on
that column, $W$ flags and its listed point--$d_b$ flags need be charged.
The table and the support bounds give at most $17\eta q$; we retain 20.
The diagonal families likewise give
\begin{equation}\label{diagonalagree}
 |\{t:g(t)\ne f(t)\}|\leq36\eta q.
\end{equation}

We supply the additive correction argument for the cross-field codomain.
Suppose $A:\F_q\to\F_r$ fails additivity on at most a fraction
$\theta<1/6$ of pairs. Put $z_x(y)=A(x+y)-A(y)$. For fixed $x$, the
additive equations at $(x+y,z)$ and $(y,x+z)$ imply $z_x(y)=z_x(z)$
except on at most a fraction $2\theta$ of independent $(y,z)$. Hence
$z_x$ has a unique mode $H(x)$ of mass at least $1-2\theta$.
The identity
\[
 z_{x+t}(y)=z_x(t+y)+z_t(y)
\]
and the three mode bounds prove additivity of $H$, since $6\theta<1$.
If $A(x)\ne H(x)$, at least $1-2\theta$ of the additive equations in
that row fail. Summing rows shows that $A$ differs from $H$ on at most
$2\theta q$ inputs.

Apply the Pexider equation to
$(a+x,b+y)$, $(a+x+y,b)$, $(a,b+y)$, and $(a+y,b)$.
Each pair is uniform for independent uniform $a,b,x,y$. Subtraction
cancels the $g,c$ terms; averaging chooses $a$ for which
$A(x)=f(a+x)-f(a)$ has additive defect at most $4r_0$.
The preceding correction supplies an additive $H$ and a constant
$\alpha$ such that $f=H+\alpha$ outside at most $8r_0q$ inputs.
Substitution into the original equation gives
\[
 g(z)-H(z)=\alpha+c(b)-H(b)
\]
outside a fraction $9r_0$ of independent pairs $(z,b)$. The largest atom
$\beta$ of $c(b)-H(b)$ has mass at least $1-9r_0$. Consequently
\begin{equation}\label{correction}
\begin{aligned}
 f&=H+\alpha &&\text{outside at most }8r_0q\text{ inputs},\\
 g&=H+\alpha+\beta &&\text{outside at most }18r_0q\text{ inputs},\\
 c&=H+\beta &&\text{outside at most }9r_0q\text{ inputs}.
\end{aligned}
\end{equation}
For the second estimate use $9r_0/(1-9r_0)\leq18r_0$.
All required inequalities, including $4r_0<1/6$ and $9r_0<1/2$,
follow from $\eta\leq10^{-7}$.
For each $i=0,1$, \eqref{fixedrow} and \eqref{correction} exclude at
most $(20\eta+27r_0)q\leq992\eta q<q$ choices of $b$. A remaining
choice gives $f(i)=H(i)+\alpha$. Equation~\eqref{protectedvalues}
therefore gives $\alpha=0$ and $H(1)=1$. Equation~\eqref{diagonalagree}
and the $f,g$ exceptions exclude at most
$(36\eta+26r_0)q\leq972\eta q<q$ choices of $t$, so $\beta=0$.
Both corrected coordinate functions are $H$.

\subsection{All slopes force one field embedding}
The affine point map equals $(H(x),H(y))$ outside at most
\[
 Z_2\leq K+26r_0q^2\leq971\eta q^2
\]
points. Mark a triple $(a,b,x)$ if $(x,ax+b)$ is a point exception or
its flag on $\ell_{a,b}$ is broken. Each exceptional affine point occurs
$q$ times. The final table family gives a total
\begin{equation}\label{marks}
 M\leq qZ_2+\eta q^3\leq972\eta q^3<q^3/4.
\end{equation}
Let
\[
 K_0=\{a\in\F_q:\text{some }s\in\F_r\text{ satisfies }
                   H(ax)=sH(x)\text{ for every }x\}.
\]
The scalar is unique because $H(1)=1$. If $a\notin K_0$, no target
line contains more than $q/2$ entries of the list
$(H(x),H(ax+b))$, $x\in\F_q$. A vertical line gives a fiber of the
nonzero additive map $H$; a nonvertical line gives a fiber of the
nonzero additive map $H(ax)-sH(x)$; the infinity line contains none.
Every fiber of a nonzero additive map has size at most $q/2$, including
when the displayed point list has repetitions. Every unmarked entry
lies on the actual image of $\ell_{a,b}$. Thus each $a\notin K_0$
and each $b$ require at least $q/2$ marks, whence
\[
 |\F_q\smallsetminus K_0|\leq2M/q^2<q/2.
\]
The set $K_0$ is a subfield. It contains 0 and 1, and scalar identities
give closure under addition and multiplication. The scalar of a nonzero
$a\in K_0$ is nonzero, since multiplication by $a$ is surjective and
$H$ is nonzero; substituting $a^{-1}x$ gives inverse closure. A proper
subfield has at most $q/2$ elements as a proper additive subgroup.
Hence $K_0=\F_q$. Its unique scalars define an injective unital field
homomorphism $\sigma$, and $H(x)=\sigma(x)H(1)=\sigma(x)$.
Equation~\eqref{correction} gives the constants $288$ and $648$.
This proves Lemma~\ref{planelemma}. All estimates charge nonnegative
flag errors, so a pointwise majorant can replace the indicator.

\section{Protection of a higher-rank frame}\label{protection}
Normalize a source frame to
$E_0=[e_0],\ldots,E_n=[e_n], I=[e_0+\cdots+e_n]$.
If it survives, normalize its target images independently to that same
coordinate frame. The first $n+1$ representatives form a basis; all
coefficients of the last representative are nonzero. Rescaling that
basis makes those coefficients 1, also in characteristic 2.
Set
\[
 H_0:X_0=0,\qquad H_i(0):X_i=0,\qquad H_i(1):X_i=X_0\quad(1\leq i\leq n).
\]
Protect $H_0$ and the $n$ hyperplanes $H_i(0)$ by their $n$ basis-point
flags each. Protect $H_i(1)$ by the flag at $I$ and its $n-1$ basis-point
flags outside $\{0,i\}$. These $2n^2+n$ original flags force the named
hyperplane images, because each contains $n$ independent named target
points.
For $j=2,\ldots,n$, let $A_j$ be the coordinate plane spanned by
$E_0,E_1,E_j$, and let $B_j$ be the vector span of the complementary
$e_k$. Its vector dimension is $n-2$; for $n=2$ it is zero. Add the
following protections, with $C_j=\{0,\ldots,n\}\smallsetminus\{0,1,j\}$:
\begin{enumerate}
\item For $S_j=[e_0+e_1+e_j]$, protect its flags with $H_k(0)$ for
$k\in C_j$ and with $H_1(1),H_j(1)$. These $n$ target hyperplanes
intersect in the specified point.
\item For $D_j:X_j=X_1$, protect its flags at $E_0,S_j$ and $E_k$ for
$k\in C_j$. Their $n$ images are independent and force $D_j$.
\item For $W_j=[e_1+e_j]$, protect its flags with $H_0,D_j$ and the
$H_k(0)$ for $k\in C_j$.
\item For $E_j^*=[e_0+e_1]$, use $H_j(0),H_1(1)$ and those $H_k(0)$.
\item For $J_j^*=[e_0+e_j]$, use $H_1(0),H_j(1)$ and those $H_k(0)$.
\end{enumerate}
The last three groups again define unique intersections of $n$ target
hyperplanes. They cost $5n$ flags per plane. Repeated flags and points
remain counted slots. The total is
\begin{equation}\label{certificate}
 M_{\rm cert}=2n^2+n+5n(n-1)=7n^2-4n\leq7n^2.
\end{equation}
Call a surviving frame protected if all these original flags are good.
For each $j$, this forces the seven points and six extended lines of
Section~\ref{plane}, with $E_1,E_j,E_0$ corresponding to $U,V,O$.
No other point-image injectivity is imposed.

\section{Transport and its unconditional load}\label{transport}
Fix a protected frame. Let $A'_j,B'_j$ be its corresponding target
subspaces. A line $\ell$ in $A_j$ extends to the source hyperplane
$\widehat\ell=\spanv(B_j,\ell)$. Define an induced point map by retaining
$\phi(p)$ when it lies in $A'_j$, and replacing it by $E_0$ otherwise.
Define the induced line image to be $\psi(\widehat\ell)\cap A'_j$ when
$\psi(\widehat\ell)$ contains $B'_j$, and the target infinity line
otherwise. In the first case the intersection is a line: a hyperplane
containing both $B'_j$ and $A'_j$ would contain their direct sum, the
whole target vector space.
Write $b_{\rm orig}$ for original broken-flag indicators, and put
\begin{equation}\label{majorant}
\begin{split}
 U_j(p,\ell)={}&b_{\rm orig}(p,\widehat\ell)
 +\sum_{k\in C_j}b_{\rm orig}(p,H_k(0))\\
 &+\sum_{k\in C_j}b_{\rm orig}(E_k,\widehat\ell).
\end{split}
\end{equation}
These are $2n-3$ original flags, counted with multiplicity. If all are
good, the middle group forces $\phi(p)\in A'_j$, the last group forces
$B'_j\subseteq\psi(\widehat\ell)$, and the first forces the induced
incidence. Thus $U_j$ pointwise majorizes induced broken flags. In
particular it pays for both possible failures of restriction. Let
$\eta_j$ be the 24-family cost evaluated on $U_j$.

For $i=1,\ldots,n$ use the original global pencil
$H_i(t):X_i=tX_0$. Call $t$ proper if $\psi(H_i(t))$ contains all $E_k$
with $k\notin\{0,i\}$ and is not $H_0$. Its equation is then
$X_i=f_i(t)X_0$; define $f_i(t)=0$ otherwise. Protection gives proper
inputs 0 and 1 with the usual values. In the induced plane $A_j$, the
proper vertical set and its extended function are exactly this global
set for $i=1$ and $f_1$. The horizontal counterparts are exactly those
for $i=j$ and $f_j$. Indeed, missing complementary containment sends
the induced line to infinity. If complementary containment holds, the
local infinity-anchor flag supplies precisely the remaining basis
direction. The intersection is the local infinity line exactly when
its span with $B'_j$ is $H_0$. The extensions by zero therefore also
agree at every nonproper input.

Now choose a uniform ordered source frame $F$, before imposing any
condition on its image. Source projective linear transformations act
transitively on frames and flags. Every certificate slot and every
original flag in \eqref{majorant}, at each fixed table parameter, is
an equivariant construction from $F$. Each such slot is consequently
uniform on the $m$ source flags and has expected error $\delta$.
Equivalently each fixed slot takes any one source flag exactly
$F_{\rm count}/m$ times. This identity precedes conditioning on survival.
The raw source formulas for $U_j$ and $\eta_j$ make sense for every
frame, even when its target image does not survive. Each normalized
table sum has mass at most 1. Hence
\[
 \mathbb E\eta_j\leq24(2n-3)\delta.
\]
For $p_x=[e_0+x_1e_1+\cdots+x_ne_n]$, define
\[
 \gamma(F)=q^{-n}\sum_{x\in\F_q^n}\sum_{i=1}^n
 b_{\rm orig}(p_x,H_i(x_i)),\qquad
 \mathcal E(F)=\gamma(F)+\sum_{j=2}^n\eta_j(F).
\]
The same equivariance gives $\mathbb E\gamma=n\delta$ and
\begin{equation}\label{expectation}
 \mathbb E\mathcal E\leq[n+24(n-1)(2n-3)]\delta\leq48n^2\delta.
\end{equation}
The certificate fails with probability at most $7n^2\delta$.
Markov's inequality and \eqref{expectation} show that the probability
of certificate failure or $\mathcal E>10^{-7}$ is at most
\[
 (7+48\cdot10^7)n^2\delta<5\cdot10^8n^2\delta\leq\rho/2.
\]
Since survival has probability $\rho>0$, some surviving protected frame
has $\mathcal E\leq10^{-7}$. If $\delta=0$, every raw nonnegative
cost is zero and the same conclusion follows directly. No conditional
uniformity or division by zero is needed.

\section{Coherence and reconstruction of the affine chart}\label{coherence}
Fix that frame. Lemma~\ref{planelemma} supplies, for each $j$, an
embedding $\sigma_j:\F_q\hookrightarrow\F_r$ with
\[
 |\{t:f_1(t)\ne\sigma_j(t)\}|\leq288\eta_jq,\qquad
 |\{t:f_j(t)\ne\sigma_j(t)\}|\leq648\eta_jq.
\]
There are at most $17\eta_jq$ nonproper inputs in either pencil.
For two charts the embeddings disagree on at most
$288(\eta_j+\eta_k)q\leq288\cdot10^{-7}q<q/2$ inputs, since both
approximate the same original function $f_1$. Distinct additive maps
can agree on at most $q/2$ inputs, by the kernel bound for their
difference. Thus all $\sigma_j$ equal one $\sigma$. For $n=2$ there
is just one chart.
Write $\eta_\Sigma=\sum_{j=2}^n\eta_j$. Use chart 2 once for coordinate
1 and chart $j$ for coordinate $j$. Summing bounds, the number of
nonproper coordinate inputs is at most
$17(\eta_2+\eta_\Sigma)q\leq34\eta_\Sigma q$, and the number of
scalar disagreements is at most
$(288\eta_2+648\eta_\Sigma)q\leq936\eta_\Sigma q$.
If all $x_i$ are proper and all $n$ coordinate flags at $p_x$ are good,
$\phi(p_x)$ is the unique intersection of
$X_i=f_i(x_i)X_0$, $1\leq i\leq n$. The union bound across the whole
affine chart gives at most $(\gamma+34\eta_\Sigma)q^n$ failures of
this product identity. After correcting the scalar values, the point
map agrees with $(x_i)\mapsto(\sigma(x_i))$ outside at most
\begin{equation}\label{affine}
 Z\leq(\gamma+970\eta_\Sigma)q^n\leq971\mathcal E q^n.
\end{equation}
This is agreement on the entire affine chart, rather than only on the
coordinate planes. Extend this embedding projectively to $G$. Its
image of $H_0$ already equals $\psi(H_0)$.

\section{Both maps are close, including at infinity}\label{infinity}
Two intersection bounds apply also to proper subfield embeddings.
A target hyperplane different from $G(H)$ contains at most $\lambda$
images of source points of $H$. A target point different from $G(p)$
lies on at most $\lambda$ images of source hyperplanes through $p$.
To see the first, normalize $G$ to coordinatewise $\sigma$. The
restriction of the wrong target equation to the $n$-dimensional
$\F_q$ vector space defining $H$ is a nonzero $\sigma$-linear
functional into $\F_r$. Its kernel has $\F_q$-dimension at most $n-1$.
For the second, the $\F_q$ coefficient space of hyperplanes through
$p$ has dimension $n$ and spans the corresponding target coefficient
space over $\F_r$. Evaluation at a target point other than $\sigma(p)$
is nonzero on that space and has the same kernel bound.
For an affine source hyperplane, a wrong target hyperplane contains
at most $q^{n-2}$ correctly embedded affine points. On its affine
parameter space, the equation is either nonconstant, with fibers of
size at most $q^{n-2}$, or a nonzero constant, with no zeros. A zero
constant would force containment of the whole target-field span $G(H)$.

Let $B$ count wrong affine source hyperplanes, and $I_{\rm bad}$ wrong
points on $H_0$. A wrong affine hyperplane has $a=q^{n-1}$ affine
points; at most $a/q$ correctly mapped ones can preserve their flags.
Charge an affine point exception on its at most $d$ source hyperplanes:
\begin{equation}\label{wronghyperplanes}
 (a-a/q)B\leq T+dZ.
\end{equation}
A wrong infinity point lies on $d-1$ affine hyperplanes, and at most
$\lambda$ correctly mapped ones can contain its image. Each wrong
affine hyperplane contains $\lambda$ infinity points. Thus
\begin{equation}\label{wronginfinity}
 (a-1)I_{\rm bad}\leq T+\lambda B.
\end{equation}
These are separate estimates and need not charge disjoint errors.
Using $d/a\leq2$, $\lambda/(a-1)=1/(q-1)$ and $a/(a-1)\leq2$ gives
\begin{equation}\label{coarse}
 D(G)=Z+B+I_{\rm bad}\leq9Z+6T/a.
\end{equation}
Since $q^n\leq N$, $T/a\leq2\delta N$, $\mathcal E\leq10^{-7}$,
and $\delta\leq1/(10^9n^2)\leq1/(4\cdot10^9)$, this yields
\begin{equation}\label{smallD}
 D(G)/N\leq9\cdot971\cdot10^{-7}+12\delta<1/1000<1/20.
\end{equation}
Both original functions are now close to the same $G$.

\section{The paired source-design gap}\label{gap}
Let $S,\mathcal B$ be the wrong point and hyperplane sets, of sizes
$s,t$, and put $D=s+t$. Let $J$ count source incidences between these
sets. A wrong point has at most $\lambda$ correctly mapped incident
hyperplanes that can contain its wrong image. There are therefore at
least $as-J$ broken flags with wrong point and correct hyperplane.
The dual bound gives at least $at-J$ with correct point and wrong
hyperplane. The two types are disjoint, so
\begin{equation}\label{paired}
 T\geq aD-2J.
\end{equation}
Let $M$ now denote the source incidence matrix. It has degree $d$;
two distinct points share $\lambda$ hyperplanes. Hence
$MM^{\mathsf T}=a\operatorname{Id}+\lambda\mathbf1\mathbf1^{\mathsf T}$.
For $n_H=|S\cap H|$ this gives
\[
 \sum_H(n_H-ds/N)^2=as(1-s/N).
\]
Center the indicator of $\mathcal B$ and use Cauchy--Schwarz to obtain
\[
 J\leq dst/N+\sqrt{as(1-s/N)t(1-t/N)}
   \leq dD^2/(4N)+\sqrt a\,D/2.
\]
Substitution in \eqref{paired} gives
\[
 T\geq\left(a-\sqrt a-\frac{dD}{2N}\right)D.
\]
Since $a\geq2$, $\sqrt a\leq3a/4$; also $d\leq2a$.
For $D\leq N/20$ the coefficient is at least
$a-3a/4-2a/40=a/5$. Equation~\eqref{smallD} therefore proves
$D(G)\leq5T/a\leq10T/d$. Undoing the independent source and target
normalizations gives exactly one map $[v]\mapsto[A\sigma(v)]$ and its
coupled hyperplane action. Counts and entry distances are preserved.
This completes the proof of Theorem~\ref{main}.

\section{Examples and limits}
For $n=2$, the complementary sums in \eqref{majorant} are empty and
there is one chart. The support proof explicitly includes $q=2$.
Every exact coupled semilinear embedding, including a proper subfield
embedding, has $T=0$, $\rho=1$, and $D=0$; invertibility over $\F_r$
ensures survival of source frames.
Constant incident point and hyperplane maps instead have $T=0$ and
$\rho=0$. More generally, point values in a proper target subspace
$W$ and hyperplane values containing $W$ can preserve every incidence
while losing every full frame. The positive survival condition detects
this rank collapse.

For a scale check, start with an exact embedding and move one source
point $p$ to the image of a different point $b$, retaining all other
entries. Exactly $\lambda$ of the $d$ hyperplanes through $p$ also
contain $b$, so exactly $a=q^{n-1}$ flags break. The distance to the
original embedding is 1, and frames avoiding $p$ survive. This tests
the normalization $T/d$, but proves neither sharp constants nor that
the small-error hypothesis holds for every such parameter choice.
The constants are not optimized and can require very large spaces
before the hypothesis admits nonzero integer errors. The theorem
asserts existence; it supplies no implemented decoder or complexity
guarantee. It concerns canonical projective spaces, not arbitrary
non-Desarguesian planes or a general building theorem.

\section{Relation to earlier work and Erd\H{o}s Problem 85}
The original question concerns eventual monotonicity of the four-cycle
minimum-degree threshold function. That question remains open here.
This manuscript supplies no determination of $f(49)$ or new extremal
bound for it. The problem statement and its discussion were read on
17 October 2026 from Bloom's problem page~\cite{problem}. The page
attributes the question to Erd\H{o}s under its keys [Er93, p.~345],
[Er94b], [Er95], and [Er96]. Those keys do not represent fresh reading
of the original articles; the bibliography endpoints were unavailable
in the saved source review.

The saved comparison read Theorem 2.3.1 and Remark 2.3.2 of
Koll\'ar--Lieblich--Olsson--Sawin~\cite{klos}, printed pages 36--38,
from the author-hosted version whose last-modified header was
18 September 2021. That theorem assumes injection, a common finite
field, and $n>3$, and measures failures of whole lines with specified
excluded subspaces. The present theorem uses supplied hyperplane data
and original point--hyperplane flags. No reduction establishing
blanket dominance in either direction is given.
The saved comparison also read the introduction, Theorem 1, Lemma 1
and Theorem 2 context, pages 1--4, of Pankov--Vetterlein~\cite{pv},
from the author-hosted version whose header was 23 December 2021.
Its exact one-way orthogonality statements over finite-dimensional
complex Hilbert spaces are not noisy finite-field two-map bounds.
The complete proofs of these external works were not re-audited.

Koll\'ar et al. explicitly credit Blum--Luby--Rubinfeld~\cite{blr} for
linearity testing. The additive correction above belongs to that
classical tradition and is included directly to handle the required
codomain and protected values. No claim is made to originate additive
correction, semilinear geometry, or projective-design mixing.
The source review also inspected only the abstract of
Cardinali--Giuzzi--Pasini~\cite{cgp}; it establishes no subsumption
comparison. Coverage of later citing work is incomplete. Further
comparison may change the priority assessment.

\section{Formal statements and review status}
The accepted package records official theorem checks, full local import
closure replay, zero holes, and only the axioms \texttt{propext},
\texttt{Classical.choice}, and \texttt{Quot.sound}. The exact principal
mapping is as follows, in namespace \texttt{Erdos85.IncidenceRepair}:
\begin{itemize}
\item Theorem~\ref{main}: \texttt{projective\_incidence\_repair}, in
\path{Openmath/IncidenceRepair/Main.lean}, asserting
\texttt{RepairClaim K L n}.
\item Corollary~\ref{dichotomy}: \texttt{repair\_dichotomy}, in
\path{Openmath/IncidenceRepair/Corollaries.lean}.
\item Corollary~\ref{fields}:
{\small\texttt{frameSurvival\_bound\_of\_no\_field\_embedding}}, in the same file.
\end{itemize}
The numerical finite-field criterion is the classical interpretation of the formal
no-embedding hypothesis. This edition reuses existing accepted evidence
for unchanged mathematical sources. Compiling its typesetting does not
check the mathematics, and this exposition has not received a new
independent statement-by-statement review.

\section{AI use}
AI systems developed the mathematical argument, formalization and earlier
source reviews, and prepared this exposition. Robert Huynh directed the
research and is the sole named author. No human mathematical review is
claimed. Formal verification concerns precise formal statements; it
does not establish novelty or public priority. This working draft
reorganizes existing work and makes no new mathematical claim.

\begin{thebibliography}{9}
\bibitem{problem} T. Bloom, \emph{Erd\H{o}s Problem 85},
\url{https://www.erdosproblems.com/85}, and discussion
\url{https://www.erdosproblems.com/forum/thread/85}.
Accessed 17 October 2026; statement page last edited 6 December 2025.
\bibitem{klos} J. Koll\'ar, M. Lieblich, M. Olsson and W. Sawin,
\emph{The Zariski topology, linear systems, and algebraic varieties},
Theorem 2.3.1 and Remark 2.3.2,
\url{https://williamsawin.com/ReconstructionBook.pdf}.
\bibitem{pv} M. Pankov and T. Vetterlein,
\emph{A geometric approach to Wigner-type theorems},
\url{https://www.flll.jku.at/sites/default/files/u24/Wigner-type-theorems.pdf}.
\bibitem{blr} M. Blum, M. Luby and R. Rubinfeld,
Self-testing/correcting with applications to numerical problems,
\emph{Journal of Computer and System Sciences} 47(3) (1993), 549--595,
\url{https://doi.org/10.1016/0022-0000(93)90044-W}.
\bibitem{cgp} I. Cardinali, L. Giuzzi and A. Pasini,
\url{https://arxiv.org/abs/2306.17079}. Abstract-only comparison in the
saved source review.
\end{thebibliography}
\end{document}
